\maketitle

\begin{abstract}
The classical approach to dynamics, following Koopman, describes how observable
functions evolve under a known transformation $T$ — but this presupposes access
to $T$.  The prior question is: given a current observation $Q$ and a future
observation $F$, what does an observer predict about $F$ given $Q$, without
assuming the dynamics?

Making prediction precise forces an answer.  The canonical answer is the regular
conditional distribution of $F$ given $Q$, a Markov kernel
$\Pi_Q : \alpha \to \mathcal{P}(\beta)$ realising $P(F \in \cdot \mid Q = q)$
across all outcomes simultaneously.  Its existence requires disintegrating a
joint distribution — which requires a $\sigma$-additive probability measure on
the observable algebra.  That measure is not assumed here; it is the object
constructed in~\cite{mahon_paper1} from coherent structured observations
satisfying collective exhaustion.

Once $\Pi_Q$ exists, composing with $Q$ gives the map
$Q_* : \Omega \to \mathcal{P}(\beta)$ sending each sample point to its
conditional law.  The central result — the predictive factorization theorem —
shows that $Q_*$ is the minimal sufficient statistic for future prediction.
Indexing over time, the kernels $\{\Pi_t\}$ compose; their composition rule is
the Chapman--Kolmogorov equation, derived from temporal coherence rather than
assumed.  From it emerge a semigroup of Markov operators $\{K_t\}$ and the
duality $\int K_t g\, d\mu = \int g\, d(\mu \star \Pi_t)$.  When the kernels
are Dirac measures, the construction recovers the classical Koopman operator
exactly.

The output is a triple $(\mathcal{P}(\beta), \nu_{Q_*}, \{K_t\})$ — predictive
state space, stationary distribution, Markov semigroup.  Setting
$Q_t = h \circ T^t$ in a measure-preserving system identifies $Q_*$ with the
delay map $x \mapsto (h(T^n x))_{n \in \mathbb{Z}}$.  The question the
construction leaves open is when $Q_*$ is injective modulo null sets — when the
predictive state space faithfully represents the underlying state space $X$.
\end{abstract}

\begin{quote}\small
\textbf{2020 Mathematics Subject Classification.}
37A30 (primary); 47D06, 60J05, 28A99 (secondary).

\textbf{Keywords.}
Koopman operator, predictive kernel, minimal predictive state, Markov semigroup,
Koopman--Perron duality, observable dynamical system, measure-preserving system.
\end{quote}

The classical approach to dynamics, following Koopman~\cite{Koopman1931},
describes how observable functions evolve: the operator $U_T f = f \circ T$
acts by precomposing with a known transformation $T$.  This presupposes access
to $T$ as a primitive.  The question this paper asks is prior: given a current
observation $Q : \Omega \to \alpha$ and a future observation $F : \Omega \to
\beta$, what does an observer predict about $F$ given $Q$?  The dynamics, if
there are any, should follow from the answer — not precede it.

Making prediction precise immediately generates a requirement.  The canonical
answer is the regular conditional distribution of $F$ given $Q$: the unique
Markov kernel $\Pi_Q : \alpha \to \mathcal{P}(\beta)$ realising
$P(F \in \cdot \mid Q = q)$ simultaneously across all outcomes.  Its existence
rests on the Rokhlin disintegration theorem, which requires a $\sigma$-additive
probability measure on the observable $\sigma$-algebra.  That measure is not an
extra assumption here.  It is what \cite{mahon_paper1} constructs: a coherent
family of finitely additive charges on a directed system of observable algebras
extends to a genuine probability measure exactly when it satisfies collective
exhaustion, the condition ruling out persistent mass on globally vanishing
events.  The predictive construction cannot begin until that measure exists;
collective exhaustion is the condition under which it does.

Once $\Pi_Q$ exists, composing with $Q$ gives the map
$Q_* := \Pi_Q(\cdot, -) \circ Q : \Omega \to \mathcal{P}(\beta)$, sending each
sample point to its conditional law.  The predictive factorization theorem shows
that every conditional expectation of a function of $F$ factors through $Q_*$:
it is the minimal sufficient statistic for future prediction.  Indexing over
time, the kernels $\{\Pi_t\}$ compose; their composition rule is the
Chapman--Kolmogorov equation, derived from temporal coherence rather than
assumed.  This yields a semigroup of Markov operators $\{K_t\}$ and the
Koopman--Perron duality connecting them with the pushforward operators
$\{\mu \mapsto \mu \star \Pi_t\}$.  When the kernels are Dirac measures, the
construction collapses to the classical Koopman picture.

The output is a triple $(\mathcal{P}(\beta), \nu_{Q_*}, \{K_t\})$.
Setting $Q_t = h \circ T^t$ in a measure-preserving system
$(X, \mathcal{B}, \mu, T)$ with $h \in L^\infty$ identifies $Q_*$ with the
delay map $\Phi_h(x) = (h(T^n x))_{n \in \mathbb{Z}}$.  The construction then
leaves open a natural question: when is $Q_*$ injective modulo null sets — when
does the predictive state space faithfully represent the underlying state space
$X$?  The final section takes this as its starting point.

The paper is self-contained.  The connection to~\cite{mahon_paper1} is
motivational; the formalism starts from a probability space
$(\Omega, \mathcal{F}, P)$ and two measurable maps:
\begin{itemize}
  \item $Q : \Omega \to \alpha$, where $(\alpha, \mathcal{A})$ is a measurable
    space;
  \item $F : \Omega \to \beta$, where $(\beta, \mathcal{B})$ is a standard
    Borel space.
\end{itemize}
The standard Borel hypothesis on $\beta$ is used once: it is the hypothesis
under which the Rokhlin disintegration theorem~\cite{Rokhlin1952} guarantees the
existence of a regular conditional distribution.  All other results hold in
greater generality.

We write $\mathcal{P}(\beta)$ for the space of probability measures on $\beta$,
equipped with its canonical (Giry) $\sigma$-algebra.  The joint pushforward
$\rho := P \circ (Q, F)^{-1}$ on $\alpha \times \beta$ has marginals
$\rho_\alpha = P \circ Q^{-1}$ and $\rho_\beta = P \circ F^{-1}$.

The regular conditional distribution of $F$ given $Q$ is the unique (up to
$P \circ Q^{-1}$-null sets) Markov kernel $\Pi_Q : \alpha \to \mathcal{P}(\beta)$
satisfying
\[
  \int_A \Pi_Q(q, B)\, d(P \circ Q^{-1})(q) \;=\; P(Q \in A,\; F \in B)
\]
for all measurable $A \subseteq \alpha$ and $B \subseteq \beta$.

\begin{definition}[Predictive kernel]
\label{def:predictive-kernel}
The \emph{predictive kernel} is $\Pi_Q$, the disintegration kernel of
$\rho = P \circ (Q, F)^{-1}$ over its first marginal:
$\Pi_Q(q, A) = P(F \in A \mid Q = q)$.
\end{definition}

The kernel is covariant under measurable maps: if $\pi : \alpha \to \alpha'$ is
measurable and $Q' = \pi \circ Q$, then $\Pi_Q(q, \cdot) = \Pi_{Q'}(\pi(q), \cdot)$
for $(P \circ Q^{-1})$-almost every $q$, by uniqueness of disintegration.

Composing $\Pi_Q$ with $Q$ gives the map sending each sample point to its
conditional law.

\begin{definition}[Minimal predictive state map]
\label{def:minimal-predictive}
The \emph{minimal predictive state map} is
\[
  Q_* : \Omega \to \mathcal{P}(\beta), \qquad Q_*(\omega) \;=\; \Pi_Q(Q(\omega), \cdot).
\]
\end{definition}

Since $\Pi_Q$ is a Markov kernel it is measurable, and $Q_*$ inherits
measurability by composition.

\begin{theorem}[Predictive factorization]
\label{thm:predictive-factorization}
For every bounded measurable $g : \beta \to \mathbb{R}$:
\[
  E[g(F) \mid \sigma(Q)] \;=\; \int g\, dQ_*(\cdot)
  \quad P\text{-a.e.}
\]
\end{theorem}
\begin{proof}
Let $h(\omega) := \int g\, dQ_*(\omega)$.  We verify the two defining properties.

\emph{Measurability.}  The map $q \mapsto \int g\, d\Pi_Q(q, \cdot)$ is
strongly measurable; since $Q_* = \Pi_Q(Q(\cdot), \cdot)$, the map $h$ is
$\sigma(Q)$-measurable.

\emph{Set-integral identity.}  For $S = Q^{-1}(A)$:
\begin{align*}
  \int_S h\, dP
  &= \int_A \int g\, d\Pi_Q(q, \cdot)\, d(P \circ Q^{-1})(q)
  \quad\text{(pushforward)} \\
  &= \int_{A \times \beta} g(f)\, d\rho(q, f)
  \quad\text{(disintegration)} \\
  &= \int_S g(F)\, dP.
\end{align*}
The conclusion follows from a.e.\ uniqueness of conditional expectation.
\end{proof}

\begin{corollary}[Predictive sufficiency]
\label{cor:sufficiency}
$E[g(F) \mid \sigma(Q)] = E[g(F) \mid \sigma(Q_*)]$ $P$-a.e.
\end{corollary}
\begin{proof}
$\sigma(Q_*) \subseteq \sigma(Q)$, and the conditional expectation given
$\sigma(Q)$ is already $\sigma(Q_*)$-measurable by the theorem.
\end{proof}

The operator $(K_{Q_*} g)(q_*) = \int g\, dq_*$ on $\mathcal{P}(\beta)$
recasts the theorem as: $E[g(F) \mid \sigma(Q)] = (K_{Q_*} g) \circ Q_*$
$P$-a.e.

\section{Dynamics and duality}

To introduce dynamics, index over a discrete time horizon.  For each
$t \in \mathbb{N}$, let $\Pi_t : \gamma \to \mathcal{P}(\gamma)$ be a Markov
kernel on the state space $\gamma$, with $\Pi_0(q, \cdot) = \delta_q$.
A family $\{\Pi_t\}$ satisfying the Chapman--Kolmogorov equation
\[
  \Pi_{t+s}(q, \cdot) \;=\; \int_\gamma \Pi_t(q', \cdot)\, d\Pi_s(q, q')
  \quad\text{(i.e., } \Pi_{t+s} = \Pi_t \circ_\kappa \Pi_s\text{)}
\]
is a \emph{Markov semigroup kernel}.  The associated \emph{Markov operators}
are
\[
  (K_t g)(q) \;:=\; \int g\, d\Pi_t(q, \cdot), \qquad g : \gamma \to \mathbb{R}.
\]

Each $K_t$ is positive, unital, and contractive on bounded measurable functions,
with $K_0 = \mathrm{id}$ (since $K_0 g(q) = \int g\, d\delta_q = g(q)$);
positivity, unitality, and contractiveness are immediate from the Markov property.

\begin{proposition}[Semigroup property]
\label{prop:semigroup}
$K_{t+s} g = K_t(K_s g)$ for all bounded measurable $g$.
\end{proposition}
\begin{proof}
$(K_{t+s} g)(q)
= \int g\, d(\Pi_t \circ_\kappa \Pi_s)(q, \cdot)
= \int \bigl(\int g\, d\Pi_t(q', \cdot)\bigr)\, d\Pi_s(q, q')
= \int (K_t g)(q')\, d\Pi_s(q, q')
= (K_t(K_s g))(q)$,
where the exchange of integrals is justified by Bochner--Fubini.
\end{proof}

The pushforward of a measure $\mu$ on $\gamma$ through the kernel is
$P_t^* \mu := \mu \star \Pi_t = \int \Pi_t(q, \cdot)\, d\mu(q)$.

\begin{theorem}[Koopman--Perron duality]
\label{thm:koopman-perron}
For every bounded measurable $g$ and finite measure $\mu$:
\[
  \int (K_t g)\, d\mu \;=\; \int g\, d(P_t^* \mu).
\]
\end{theorem}
\begin{proof}
$\int g\, d(P_t^* \mu)
= \int g\, d(\mu \star \Pi_t)
= \int_\gamma \int g\, d\Pi_t(q, \cdot)\, d\mu(q)
= \int (K_t g)\, d\mu$,
by Bochner--Fubini.
\end{proof}

When $\Pi_t(q, \cdot) = \delta_{\varphi_t(q)}$ for a measurable map
$\varphi_t : \gamma \to \gamma$, the stochastic picture collapses.

\begin{proposition}[Deterministic limit]
\label{prop:deterministic}
If $\Pi_t(q, \cdot) = \delta_{\varphi_t(q)}$, then $K_t g = g \circ \varphi_t$
and $\varphi_{t+s} = \varphi_t \circ \varphi_s$.
\end{proposition}
\begin{proof}
$(K_t g)(q) = \int g\, d\delta_{\varphi_t(q)} = g(\varphi_t(q))$.
The flow law follows from
$\delta_{\varphi_{t+s}(q)} = (\Pi_t \circ_\kappa \Pi_s)(q, \cdot)
= \delta_{\varphi_t(\varphi_s(q))}$
and injectivity of the Dirac embedding.
\end{proof}

In the deterministic case, $K_t g = g \circ \varphi_t = U_T^t g$: the Markov
operator specialises to the classical Koopman operator.  The
Koopman--Perron duality $\int K_t g\, d\mu = \int g\, d(P_t^* \mu)$ becomes
the change-of-variables formula $\int g \circ T^t\, d\mu = \int g\, d(T^t_* \mu)$.

The full structure produced by the construction is:

\begin{definition}[Observable dynamical system]
\label{def:obs-dyn-sys}
An \emph{observable dynamical system} is a triple
$(\mathcal{P}(\beta),\, \nu_{Q_*},\, \{K_t\}_{t \in \mathbb{N}})$
where $\mathcal{P}(\beta)$ is the predictive state space (the image of $Q_*$),
$\nu_{Q_*} = P \circ Q_*^{-1}$ is the pushforward of $P$ along $Q_*$, and
$\{K_t\}$ is the Markov semigroup on $\mathcal{P}(\beta)$.
\end{definition}

Setting $Q_t = h \circ T^t$ in a measure-preserving system
$(X, \mathcal{B}, \mu, T)$ with $h \in L^\infty$ specialises the construction.
Under this choice, $Q_*(\omega) = (h(T^n \omega))_{n \in \mathbb{Z}}$ is the
delay map $\Phi_h$, and the observable dynamical system is the triple
$(\Phi_h(X), \mu \circ \Phi_h^{-1}, \{K_t\})$.

The question the construction leaves open is whether $\Phi_h$ is injective
modulo $\mu$-null sets: whether the predictive state space $\Phi_h(X)$ faithfully
represents $X$.  Injectivity modulo null sets is equivalent to the observable
$\sigma$-algebra $\mathcal{O}_h = \sigma(\{h \circ T^n : n \in \mathbb{Z}\})$
equalling $\mathcal{B}(X)$ modulo $\mu$.  The Koopman--Perron duality translates
this into a function-space question: $\mathcal{O}_h = \mathcal{B}(X)$ modulo
$\mu$ if and only if the algebra generated by $\{h \circ T^n\}$ is dense in
$L^2(X, \mu)$.  The key is that the $L^2$-closure of a function algebra equals
the $L^2$-space of the $\sigma$-algebra it generates — a density bridge between
$\sigma$-algebra generation and $L^2$ approximation.

The cyclic vector condition
$\overline{\mathrm{span}}\{U_T^n h : n \in \mathbb{Z}\} = L^2(X, \mu)$ is
sufficient but strictly stronger: reconstruction requires only that the
\emph{algebra} generated by the delay orbit be $L^2$-dense, not the linear
span.  They coincide only when $U_T$ has simple spectrum.

When reconstruction holds, $\Phi_h(X)$ is the full state space $X$ up to
measure-theoretic identification.  The compact Stone space
of~\cite{mahon_paper1} — introduced as a technical device for
$\sigma$-additivity — is, under the reconstruction condition, identified with
$X$ itself.  The construction begins with a compact completion of the sample
space and ends by recognising it as the state space.

\bibliographystyle{plain}
\bibliography{references}
